Table~\ref{tab:tdo:refine} lists the Orbiter commands for decomposition refinement.
\orbitertableofcommands{tdo_refinement}

\bigskip

Suppose we want to study projective planes of order $16.$ 
It is a linear space with $16^2+16+1=273$ points and equally many lines. 
Each point lies on $17$ lines and each line contains $17$ points. 
Any two points lie on exactly one line 
and any two lines intersect in exactly one point.

\bigskip

We decide to 
study maximal arcs of degree $4$ in this plane 
(the degree has to divide the order of the plane).
A maximal arc of degree $d$ is a set of points 
so that each line intersects in either $d$ or zero points.
A line which intersects in $d$ points is called a secant.  
A line which intersects in no point is called an external line.
The command 

\sourcecode{max_arc_16_4_start}

creates a decomposition stack for the parameters of the arc and writes the file \verb'max_arc_q16_r4.stack' 
\orbiteroutput{GRAPHICS/LATEX/max_arc_1}
This is a point-tactical decomposition with 2 point-classes and 2 block-classes. 
The point classes are associated with the rows. 
The block-classes are associated with the columns. 
The first row and column indicates the size of the classes.
The entries $a_{ij}$ count the number of blocks in the column class $j$ 
that are incident with a given point in the $i$th row  class.
The fuse information at the bottom (\verb'1 1') 
is a partition of the row classes which indicates the ancestor decomposition which was column tactical.
The next step is to convert the stack file to a tdo file.
The command

\sourcecode{max_arc_16_4_convert_stack_tdo}

does that. 
It creates the file \verb'max_arc_q16_r4.tdo'. It also 
prints the decomposition stack:
\orbiteroutput{GRAPHICS/LATEX/max_arc_2}
Next, we can compute all coarsest column-tactical refinements of the decomposition. 
To this end, the command 

\sourcecode{max_arc_16_4_refine}

is used. 
Because the incidence structure is a projective plane, 
the dual is a linear space also. 
Hence the option \verb'-dual_is_linear_space' can be used, which is helpful to reduce possibilities. 
As it turns out, there is exactly one refinement, and it is tactical. 
The file \verb'max_arc_q16_r4r.tdo' is produced. 
Note the added letter \verb'r' at the end of the file name (r for refinement).
We can use the following command to display the decomposition stack in the file:

\sourcecode{max_arc_16_4r_print}

This produces the following output:
\orbiteroutput{GRAPHICS/LATEX/max_arc_3}

